\section{Мощности множеств}

Во втором доказательстве существования алгебраического замыкания нужна оценка мощности алгебраических расширений.
Соберём необходимые факты. Аксиома выбора используется в \rf{6.4e}--\rf{6.4h}.

\begin{definition}[6.4b]
Множества $X$ и $Y$ \emph{равномощны}, $|X|=|Y|$, если существует биекция $X\to Y$; $|X|\le|Y|$, если существует
инъекция $X\to Y$; $|X|<|Y|$, если $|X|\le|Y|$, но $|X|\ne|Y|$. Множество $X$ \emph{счётно}, если оно равномощно
$\mathbb N=\{0,1,2,\dots\}$ (пишут $|X|=\aleph_0$), и \emph{не более чем счётно}, если $|X|\le\aleph_0$ (т.\,е.\ конечно или
счётно). Сумма и произведение мощностей: $|X|+|Y|=|X\sqcup Y|$ (дизъюнктное объединение) и
$|X|\cdot|Y|=|X\times Y|$.
\end{definition}

\begin{theorem}[6.4c: Кантор--Шрёдер--Бернштейн]
Если $|X|\le|Y|$ и $|Y|\le|X|$, то $|X|=|Y|$.
\end{theorem}

\begin{proof}
Пусть $f\colon X\to Y$ и $g\colon Y\to X$ инъективны. Положим $C_0=X\setminus g(Y)$, $C_{n+1}=g(f(C_n))$ и
$C=\bigcup_{n\ge0}C_n$. Если $x\notin C$, то $x\notin C_0$, т.\,е.\ $x\in g(Y)$. Определим $h\colon X\to Y$ так:
$h(x)=f(x)$ при $x\in C$ и $h(x)=g^{-1}(x)$ при $x\notin C$.

\emph{Инъективность.} На $C$ и на $X\setminus C$ отображение $h$ инъективно. Если $f(x)=g^{-1}(x')$ при $x\in C_n$ и
$x'\notin C$, то $x'=g(f(x))\in C_{n+1}\subset C$~--- противоречие.

\emph{Сюръективность.} Пусть $y\in Y$. Если $g(y)\notin C$, то $h(g(y))=y$. Если $g(y)\in C$, то $g(y)\notin C_0$,
значит, $g(y)\in C_{n+1}=g(f(C_n))$ для некоторого $n$, т.\,е.\ $y=f(x)$ с $x\in C_n$ ($g$ инъективно), и $h(x)=y$.
\end{proof}

\begin{theorem}[6.4d: Кантор]
Для любого множества $X$ имеем $|X|<|\mathcal P(X)|$, где $\mathcal P(X)$~--- множество всех подмножеств $X$.
\end{theorem}

\begin{proof}
Отображение $x\mapsto\{x\}$ инъективно. Пусть $s\colon X\to\mathcal P(X)$~--- любое отображение и
$D=\{x\in X\mid x\notin s(x)\}$. Если $D=s(d)$, то $d\in D\iff d\notin s(d)=D$~--- противоречие. Значит, $s$ не
сюръективно, и биекции $X\to\mathcal P(X)$ нет.
\end{proof}

Таким образом, для любого множества есть множество строго большей мощности.

\begin{lemma}[6.4e: сравнимость мощностей]
Для любых множеств $X$, $Y$ либо $|X|\le|Y|$, либо $|Y|\le|X|$.
\end{lemma}

\begin{proof}
Рассмотрим множество $\mathcal I$ инъекций $u\colon A\to Y$, определённых на подмножествах $A\subset X$, с порядком
<<продолжения>>: $u\le u'$, если $u'$ продолжает $u$. Оно непусто (пустое отображение), а у цепи есть верхняя грань~---
<<объединение>> отображений цепи (оно инъективно, так как любые две точки лежат в области определения одного
отображения цепи). По лемме Цорна есть максимальная инъекция $u\colon A\to Y$. Если бы нашлись $x\in X\setminus A$ и
$y\in Y\setminus u(A)$, то $u$ продолжалась бы правилом $x\mapsto y$. Значит, либо $A=X$, и тогда $|X|\le|Y|$, либо
$u(A)=Y$, и тогда $u^{-1}\colon Y\to A\subset X$~--- инъекция, $|Y|\le|X|$.
\end{proof}

В частности, в каждом бесконечном множестве $X$ есть счётное подмножество: по \rf{6.4e} либо $\mathbb N$ вкладывается в $X$
(и образ счётен), либо $|X|\le\aleph_0$, и тогда $X$ равномощно бесконечному подмножеству $\mathbb N$, а оно счётно
(перечисляется по возрастанию).

\begin{theorem}[6.4f]
Если множество $X$ бесконечно, то $|X\times X|=|X|$.
\end{theorem}

\begin{proof}
\emph{Счётный случай.} Отображение $c(m,n)=\frac{(m+n)(m+n+1)}2+n$~--- биекция $\mathbb N\times\mathbb N\to\mathbb N$
(\emph{нумерация Кантора}): пары перебираются по диагоналям $m+n=d$, на диагонали $d$ лежит $d+1$ пара, и им достаются
номера от $\frac{d(d+1)}2$ до $\frac{d(d+1)}2+d$.

\begin{center}
\begin{tikzpicture}[scale=0.82,>=stealth]
\draw[->,gray] (-0.4,0) -- (5.1,0) node[right]{$m$};
\draw[->,gray] (0,-0.4) -- (0,5.1) node[above]{$n$};
\foreach \d in {1,...,4} \draw[->,cnew,thick,shorten <=4pt,shorten >=4pt] (\d,0) -- (0,\d);
\foreach \d in {0,...,3} \draw[->,gray,dashed,shorten <=4pt,shorten >=4pt] (0,\d) -- (\d+1,0);
\foreach \d in {0,...,4} \foreach \n in {0,...,\d} {
  \pgfmathtruncatemacro{\m}{\d-\n}
  \pgfmathtruncatemacro{\c}{\d*(\d+1)/2+\n}
  \fill (\m,\n) circle (1.6pt);
  \node[font=\scriptsize,above right=-1.5pt] at (\m,\n) {$\c$};
}
\end{tikzpicture}
\end{center}
\figcap{Нумерация Кантора: точка $(m,n)$ получает номер $c(m,n)$; номера идут по диагоналям (красные стрелки) с
переходом на следующую диагональ (пунктир).}

\emph{Общий случай.} Пусть $\mathcal S$~--- множество пар $(A,f)$, где $A\subset X$ бесконечно, а
$f\colon A\to A\times A$~--- биекция; положим $(A,f)\le(B,g)$, если $A\subset B$ и $g|_A=f$.

\emph{1) $\mathcal S$ непусто:} в $X$ есть счётное подмножество $A_0$, а для него биекция $A_0\to A_0\times A_0$ даётся
нумерацией Кантора.

\emph{2) У каждой цепи есть верхняя грань.} Для пустой цепи годится любой элемент $\mathcal S$ (он есть по 1)). Для
непустой цепи $\{(A_i,f_i)\}$ возьмём $(A,f)$, где $A=\bigcup A_i$ и $f=\bigcup f_i$ (определено
корректно, так как $f_i$ согласованы). Отображение $f$ инъективно: любые два элемента $A$ лежат в одном $A_k$ (цепь
линейно упорядочена), а $f_k$ инъективно. Оно сюръективно: если $a\in A_i$, $b\in A_j$ и $A_k$~--- большее из
$A_i,A_j$, то $(a,b)\in A_k\times A_k=f_k(A_k)$.

По лемме Цорна есть максимальный элемент $(A,f)$; положим $\kappa=|A|$, так что $\kappa\cdot\kappa=\kappa$.

\emph{3) $3\kappa=\kappa$ (и тем более $2\kappa=\kappa$).} Возьмём различные $a_0,a_1,a_2\in A$; отображение
$A\times\{0,1,2\}\to A\times A$, $(a,i)\mapsto(a,a_i)$, инъективно, поэтому $3\kappa\le\kappa\cdot\kappa=\kappa$.
Неравенства $\kappa\le2\kappa\le3\kappa$ очевидны; остаётся \rf{6.4c}.

\emph{4) $|X\setminus A|<\kappa$.} Иначе (\rf{6.4e}) $|X\setminus A|\ge\kappa$, и найдётся $C\subset X\setminus A$ с
$|C|=\kappa$. Тогда
\[
(A\sqcup C)\times(A\sqcup C)=(A\times A)\sqcup\bigl((A\times C)\sqcup(C\times A)\sqcup(C\times C)\bigr),
\]
и каждое из трёх последних множеств имеет мощность $\kappa\cdot\kappa=\kappa$, а их объединение~--- мощность
$3\kappa=\kappa=|C|$. Значит, есть биекция $g\colon C\to(A\times C)\sqcup(C\times A)\sqcup(C\times C)$, и
$f\cup g\colon A\sqcup C\to(A\sqcup C)\times(A\sqcup C)$~--- биекция, продолжающая $f$. Это противоречит
максимальности $(A,f)$.

\emph{5) Вывод.} $X=A\sqcup(X\setminus A)$, поэтому $\kappa\le|X|\le\kappa+\kappa=\kappa$, т.\,е.\ $|X|=\kappa$ (\rf{6.4c}), и
$|X\times X|=\kappa\cdot\kappa=\kappa=|X|$.
\end{proof}

\begin{corollary}[6.4g]
Пусть $X$ бесконечно. Тогда
\begin{enumerate}[label=(\alph*)]
\item $|X\sqcup X|=|X|$ и $|X^n|=|X|$ для всех $n\ge1$;
\item $\bigl|\bigsqcup_{n\ge1}X^n\bigr|=|X|$;
\item если $|T|\le|X|$ и $(S_t)_{t\in T}$~--- семейство конечных множеств, то $\bigl|\bigcup_{t\in T}S_t\bigr|\le|X|$.
\end{enumerate}
\end{corollary}

\begin{proof}
(a) $X\sqcup X$ вкладывается в $X\times X$ (возьмём две точки $x_0\ne x_1$ и пошлём $x$ из $i$-й копии в $(x,x_i)$),
поэтому $|X|\le|X\sqcup X|\le|X\times X|=|X|$ (\rf{6.4f}), и по \rf{6.4c}
$|X\sqcup X|=|X|$. Далее индукция: $|X^n|=|X^{n-1}\times X|=|X\times X|=|X|$.

(b) Выберем биекции $X^n\to X$; они дают инъекцию $\bigsqcup_nX^n\to\mathbb N\times X$, а $\mathbb N\times X$
вкладывается в $X\times X$ (в $X$ есть счётное подмножество). Обратное неравенство очевидно.

(c) Выберем инъекцию $T\to X$, для каждого $t$~--- инъекцию $e_t\colon S_t\to\mathbb N$, а для каждого
$s\in\bigcup S_t$~--- некоторое $t(s)$ с $s\in S_{t(s)}$. Отображение $s\mapsto\bigl(t(s),e_{t(s)}(s)\bigr)$~---
инъекция $\bigcup S_t\to T\times\mathbb N$, а $|T\times\mathbb N|\le|X\times X|=|X|$.
\end{proof}

\begin{proposition}[6.4h]
Пусть $E/F$~--- алгебраическое расширение.
\begin{enumerate}[label=(\alph*)]
\item Если $F$ бесконечно, то $|E|=|F|$.
\item Если $F$ конечно, то $E$ конечно или счётно.
\end{enumerate}
\end{proposition}

\begin{proof}
Многочлен степени меньше $n$ задаётся набором из $n$ коэффициентов, поэтому
$|F[X]|\le\bigl|\bigsqcup_{n\ge1}F^n\bigr|$. Каждый $\alpha\in E$~--- корень своего минимального многочлена
$\mu_\alpha\in F[X]$, а у каждого ненулевого многочлена конечное число корней в $E$. Значит, $E$~--- объединение
семейства конечных множеств (корней в $E$ ненулевых многочленов из $F[X]$), занумерованного множеством мощности не
больше $|F[X]|$.

(a) По \rf{6.4g}(b) $|F[X]|\le|F|$, а по \rf{6.4g}(c) $|E|\le|F|$. Так как $F\subset E$, то $|F|\le|E|$, и по \rf{6.4c} $|E|=|F|$.

(b) Здесь $F[X]$~--- объединение счётного числа конечных множеств, т.\,е.\ не более чем счётно (\rf{6.4g}(c) при
$X=\mathbb N$), и $E$~--- объединение не
более чем счётного семейства конечных множеств, т.\,е.\ не более чем счётно.
\end{proof}

\begin{remark}[6.4i]
\begin{enumerate}[label=(\alph*),leftmargin=*]
\item \emph{Где нужен выбор.} Утверждение \rf{6.4f} для всех бесконечных $X$ равносильно аксиоме выбора; выбор используется и
в \rf{6.4g} (выбор биекций и инъекций). Для самих полей, встречающихся на практике, равенства $|F\times F|=|F|$ и
$|F[X]|=|F|$ получаются без выбора: счётные поля ($\Q$, числовые поля, $\Q(T)$, $\F_p(T)$) перечисляются явно, а для
$\R$ явную биекцию $\R\times\R\to\R$ даёт равенство $|\R|=|\{0,1\}^{\mathbb N}|$ и <<перемешивание>> двух
последовательностей из нулей и единиц в одну: первая~--- на чётных местах, вторая~--- на нечётных. Но в \rf{6.4h} выбор
нужен даже при $F=\Q$: в \rf{6.4g}(c) выбирается нумерация корней каждого многочлена в $E$. Без выбора алгебраическое
расширение $\Q$ может оказаться несчётным: таково алгебраическое замыкание $\Q$ без нетривиальных автоморфизмов из
\rf{6.4a}. Будь оно счётным, его нумерация позволила бы провести счётный случай \rf{6.6} без выбора и получить
изоморфизм с $\Qal$, у которого есть комплексное сопряжение.
\item \emph{Перечислением не обойтись.} Доказательство \rf{1.33} исчерпывает множество алгебраических чисел возрастающей
последовательностью конечных множеств. Для несчётного $F$ так нельзя: (при аксиоме выбора) несчётное множество не является
объединением последовательности конечных множеств. Для произвольного $F$ нужна арифметика мощностей \rf{6.4f}--\rf{6.4g}.
\item \emph{Трансцендентные числа} (Кантор). Множество алгебраических чисел счётно (\rf{1.33}; это также \rf{6.4h}(a) для
$F=\Q$), а $\CC$ несчётно. Поэтому трансцендентные числа существуют, и <<почти все>> комплексные числа трансцендентны.
\item При трансцендентных расширениях мощность может вырасти: $|\CC|>|\Q|$. Для алгебраического расширения $\CC/\R$, как
и должно быть, $|\CC|=|\R|$.
\end{enumerate}
\end{remark}
